\section{Lecture 3}
\label{sec:lecture-3}
\setcounter{equation}{0}
\subsection*{Torsion forms and Bott-Chern secondary class}
\subsubsection*{Chern connection}
...

\begin{theorem}
    For any holomorphic vector bundle $E$ on a complex manifold $X$ with 
    a Hermitian metric $h^E$, there exists a unique connection 
    $\nabla^E$ such that
    \begin{enumerate}
        \item $\nabla^E$ is compatible with $h^E$, i.e. for any 
        $U\subset X$ open, $s_1,s_2\in C^\infty(U,E)$, 
        \[
        d\langle s_1,s_2\rangle_{h^E} = 
        \langle\nabla^E s_1,s_2\rangle_{h^E} + 
        \langle s_1,\nabla^E s_2\rangle_{h^E};
        \]
        \item $\nabla^E$ is compatible with the holomorphic structure of 
        $E$, i.e. $\nabla^{E''}=\bar{\partial}^E$.
    \end{enumerate}
    This connection is called the Chern connection on $(E,h^E)$.
\end{theorem}

For now on, we always denote by $\nabla^E$ the Chern connection. 

\subsubsection*{Torsion forms for Chern complex}
\textbf{Settings:}
Let 
\[
(E,v):0\to E_1\overset{v}{\to}\cdots\overset{v}{\to}E_m\to0
\]
be a holomorphic chain complex of finite dimensional holomorphic vector 
bundles on a connected complex manifold $X$, i.e. $E_j$ are holomorphic 
vector bundles and $v:E_j\to E_{j+1}$ are $\mC$-linear holomorphic map. 
\textcolor{blue}{(If $s$ is a local holomorphic section of $E_j$ on $U$, 
then $(vs)_x:=v_xs_x$ is a local holomorphic section of $E_{j+1}$ on $U$ 
$\Longleftrightarrow$ $\bar{\partial}^{E}v+v\bar{\partial}^{E} = 0$).}
We say it is a chain complex means: 
\[
v^2=0\Longleftrightarrow(v+\bar{\partial}^{E})^2=0\text{ on }
\Omega^{\cdot,\cdot}(X,E).
\]

Set 
\[
E_+=\bigoplus_{j\text{ even}}E_j,\;
E_-= \bigoplus_{j\text{ odd}}E_j,\;
E=E_+\oplus E_-.
\]
Let $h^E=\bigoplus\limits_j h^{E_j}$ be a Hermitian metric on $E$. 
So that $E_i\perp E_j$ if $i\neq j$. 

Let $v^*$ be the adjoint of $v$ w.r.t. $h^E$. Set $V=v+v^*$. 
Then $\nabla^E+V$ is a superconnection on $E$.

Let 
\begin{gather*}
    P^X:=\bigoplus_p\Omega^{p,p}(X), \\
    P^{X,0}:=\{\omega\in P^X\mid\exists\,\eta,\eta'\in\Omega(X)
    \;s.t.\;\omega=\partial\eta+\bar{\partial}\eta'\}
\end{gather*}

\begin{proposition}
    Set 
    \[
    A=\bigoplus_{p,q,j}\Omega^{p,q}\big(X,\Hom(E_j,E_{j+p-q})\big),
    \]
    then $A$ is a super-algebra. Moreover, if $\eta\in A$, then 
    $\Tr_s[\eta]\in P^X$.
\end{proposition}
\begin{proof}
    If $\eta,\eta'\in A$, then $\eta\,\hat{\otimes}\,\eta'\in A$. 
    Since $\Tr_s$ vanishes on $\Hom(E_j,E_k)_{j\neq k}$, thus 
    $\Tr_s[\eta]\in P^X$.
\end{proof}

Let 
\[
A_u=\nabla^E+\sqrt{u}V,\quad u\geq0.
\]
be a superconnection on $E$. Let $N$ be a number operator on $E$,
i.e., 
\[
N|_{E_j}=j\Id_{E_j}.
\] 
\begin{theorem}
    For any $u\geq0$, $\Tr_s[\exp(-A_u^2)],\,\Tr_s[N\exp(-A_u^2)]$
    are in $P^X$. The form $\Tr_s[\exp(-A_u^2)]$ is closed. 
    For $u>0$, 
    \begin{gather}
    \pd{}{u}\Tr_s[\exp(-A_u^2)] = 
    -\frac{1}{2u}(\p+\bp)\Tr_s[\sqrt{u}V\exp(-A_u^2)],\label{eq:3.1} \\
    \Tr_s[\sqrt{u}V\exp(-A_u^2)] = 
    (\p-\bp)\str[N\exp(-A_u^2)].\label{eq:3.2}
    \end{gather}
    In particular, 
    \begin{equation}\label{eq:3.3}
        \pd{}{u}\Tr_s[\exp(-A_u^2)] = 
        -\frac{1}{u}\bp\p\Tr_s[N\exp(-A_u^2)].
    \end{equation}
\end{theorem}
\begin{proof}
    Step $1$: Show that $A_u^2\in A$.
    
    Since 
    \begin{equation*}
    A_u^2=(\nabla^E)^2+uV^2+\sqrt{u}[\nabla^E,V],
    \end{equation*}
    where
    \begin{gather*}
        (\nabla^E)^2\in\Omega^{1,1}(X,\End(E)), \\
        V^2=(v+v^*)^2=vv^*+v^*v\in C^\infty(X,\End(E)), \\
    \end{gather*}
    and
    \begin{equation*}
        [\nabla^E,V]=[\nabla'',v] + [\nabla',v] + 
        [\nabla'',v^*] + [\nabla',v^*]
    \end{equation*}
    By the definition of holomorphic chain complex, 
    \begin{equation*}
        [\nabla'',v] = \bp v + v\bp = 0.
    \end{equation*}
    For $u\in T^{1,0}X$, 
    \begin{align*}
        \langle([\nabla'_u,v^*]\sigma_1,\sigma_2)\rangle &= 
        \langle(\nabla'_u v^*-v^* \nabla'_u)\sigma_1,\sigma_2\rangle
        =\langle\nabla'_u(v^*\sigma_1),\sigma_2\rangle 
        -\langle\nabla'_u\sigma_1,v\sigma_2\rangle \\
        &=u\langle v^*\sigma_1,\sigma_2\rangle 
        -\langle v^*\sigma_1,\nabla''_{\bar{u}}\sigma_2\rangle 
        -u\langle\sigma_1,v\sigma_2\rangle 
        +\langle\sigma_1,\nabla''_{\bar{u}}(v\sigma_2)\rangle \\
        &=\langle\sigma_1,
        (\nabla''_{\bar{u}}v-v\nabla''_{\bar{u}})\sigma_2\rangle \\
        &=0.
    \end{align*}
    \textcolor{gray}{
    The last equality holds since $\nabla''=\bp$ and notice that 
    $v$ acts on $\Omega^*(X,E)$ by 
    \begin{equation*}
        v(\alpha\otimes\sigma) = (-1)^{\deg\alpha}\alpha\otimes(v\sigma).
    \end{equation*}
    locally $\bp\sigma_2=\alpha^i\otimes e_i$, $\deg\alpha^i=1$, then 
    \begin{align*}
        i_{\bar{u}}v(\bp\sigma_2) 
        &= i_{\bar{u}}v(\alpha^i\otimes e_i) 
        = -i_{\bar{u}}\big(\alpha^i\otimes(v e_i)\big) \\
        &= -\alpha^i(\bar{u})\otimes(v e_i)
        = -v\big(\alpha^i(\bar{u})\otimes e_i\big) \\
        &= -v(i_{\bar{u}}\bp\sigma_2)
        = -v(\nabla''_{\bar{u}}\sigma_2).
    \end{align*}
    thus $(\bp v + v\bp)\sigma_2=0$ implies 
    $(\nabla''_{\bar{u}}v - v\nabla''_{\bar{u}})\sigma_2=0$.
    }

    So 
    \begin{equation*}
        [\nabla^E,V]=[\nabla',v] + [\nabla'',v^*]
    \end{equation*}
    where
    \begin{equation*}
        [\nabla',v]\in \Omega^{1,0}(X,\End(E_j,E_{j+1})),\quad 
        [\nabla'',v^*]\in \Omega^{0,1}(X,\End(E_j,E_{j-1}))
    \end{equation*}
    Thus $A_u^2\in A$, which deduces $\exp(-A_u^2)\in A$. With $N\in A$, 
    we conclude that $\Tr_s[\exp(-A_u^2)],\,\Tr_s[N\exp(-A_u^2)]\in P^X$.

    Step $2$: Verify the formulas via Duhamel's formula.
    \begin{lemma}
        For $H_u\in \Big(\bigwedge(T^*X)\otimes\End(E)\Big)^+$ 
        \begin{equation*}
            \pd{}{u}\exp(-tH_u) = 
            -\int_0^t\me^{-(t-s)H_u}(\pd{}{u}H_u)\me^{-sH_u}\md s.
        \end{equation*}
    \end{lemma}

    \textcolor{gray}{
    \begin{proof}[Proof of lemma]
        Recall that $\exp(-tH)$ is the unique solution of the ODE
        \begin{equation*}
            \left\{
            \begin{aligned}
                &\Big(\pd{}{t}+H\Big)u_t=0, \\
                &u_0=\Id_E.
            \end{aligned}
            \right.
        \end{equation*}
        In order to verify the formula, we only need to check that
        \begin{equation*}
            Q_t:=\pd{}{u}\exp(-tH_u)
            +\int_0^t\me^{-(t-s)H_u}(\pd{}{u}H_u)\me^{-sH_u}\md s
        \end{equation*}
        satisfies the ODE 
        \begin{equation*}
            \left\{
            \begin{aligned}
                &\Big(\pd{}{t}+H\Big)Q_t=0, \\
                &Q_0=0.
            \end{aligned}
            \right.
        \end{equation*}
        Indeed,
        \begin{align*}
            \pd{}{t}Q_t &= \pd{}{u}\pd{}{t}\exp(-tH_u)
            +\me^{-(t-s)H_u}(\pd{}{u}H_u)\me^{-sH_u}\Big|_{s=t} \\
            &\quad+\int_0^t\pd{}{t}\Big(\me^{-(t-s)H_u}(\pd{}{u}H_u)
            \me^{-sH_u}\Big)\md s \\
            &= -\pd{}{u}(H_u\exp(-tH_u)) + (\pd{}{u}H_u)\me^{-tH_u} \\
            &\quad+\int_0^t -H_u\me^{-(t-s)H_u}(\pd{}{u}H_u)\me^{-sH_u}\md s \\
            &= -H_u\pd{}{u}\exp(-tH_u)
            -\int_0^t \me^{-(t-s)H_u}(\pd{}{u}H_u)\me^{-sH_u}\md s \\
            &= -H_u Q_t.
        \end{align*}
        And $Q_0=\pd{}{u}\Id_E+0=0$. Thus $Q_t=0, \forall t$.
    \end{proof}
    }

    Now for formula \eqref{eq:3.1}, by lemma we have
    \begin{align*}
        &\pd{}{u}\str\left[\exp(-A_u^2)\right] = 
        -\str\left[\int_{0}^{1}\me^{-(1-s)A_u^2}(\pd{}{u}A_u^2)
        \me^{-sA_u^2}\md s\right] \\
        =&-\int_{0}^{1}\str\left[(\pd{}{u}A_u^2)\me^{-(1-s)A_u^2}
        \me^{-sA_u^2}\right]\md s \quad\textcolor{gray}{\str[A,B]=0} \\
        =&-\str\left[\pd{}{u}A_u^2,\me^{-A_u^2}\right] \\
        =&-\str\left[\left[A_u,\frac{1}{2\sqrt{u}}V\right],
        \me^{-A_u^2}\right] \quad
        \textcolor{gray}{\pd{}{u}A_u^2=\left[A_u,\pd{}{u}A_u\right],}\,
        \textcolor{gray}{\pd{}{u}A_u = \frac{1}{2\sqrt{u}}V} \\
        =&-\str\left[A_u,\frac{1}{2\sqrt{u}}V\me^{-A_u^2}\right]\quad 
        \textcolor{gray}{[A_u,A_u^2]=0} \\
        =&-\md\str\left[\frac{1}{2\sqrt{u}}V\me^{-A_u^2}\right] \\
        =&-\frac{1}{2u}(\p+\bp)\Tr_s[\sqrt{u}V\exp(-A_u^2)]
    \end{align*}

    To verify formula \eqref{eq:3.2}, we introduce a parameter $b$ and study 
    \begin{equation*}
        \Tr_s\left[\exp\left(-A_u^2+bN\right)\right].
    \end{equation*}
    Since 
    \begin{equation*}
        \md\str\left[\exp\left(-A_u^2+bN\right)\right]
        =\str\left[\left[A_u,\exp\left(-A_u^2+bN\right)\right]\right]
    \end{equation*}
    To calculate the right hand side, we set 
    \begin{equation*}
        Q_t:=\left[A_u,\exp\left(-A_u^2+bN\right)\right] + 
        \int_0^t\me^{-(t-s)(A_u^2-bN)}[A_u,A_u^2-bN]
        \me^{-s(A_u^2-bN)}\md s.
    \end{equation*}
    similarly we have
    \begin{align*}
        \pd{}{t}Q_t &= \left[A_u,-(A_u^2-bN)\exp\left(-t(A_u^2-bN)\right)
        \right] + \left[A_u,A_u^2-bN\right]\me^{-t(A_u^2-bN)} \\
        &\quad-\int_0^t(A_u^2-bN)\Big(\me^{-(t-s)(A_u^2-bN)}
        [A_u,A_u^2-bN]\me^{-s(A_u^2-bN)}\Big)\md s \\
        &= -(A_u^2-bN)\left[A_u,\exp\left(-t(A_u^2-bN)\right)\right] \\
        &\quad-(A_u^2-bN)\int_0^t\me^{-(t-s)(A_u^2-bN)}
        [A_u,A_u^2-bN]\me^{-s(A_u^2-bN)}\md s \\
        &= -(A_u^2-bN)Q_t.
    \end{align*}
    And $Q_0=[A_u,\Id_E] +0=0$. Thus $Q_t=0, \forall t$. Thus 
    \begin{align*}
        &\md\str\left[\exp\left(-A_u^2+bN\right)\right] \\
        =& -\int_0^1\str\left[\me^{-(1-s)(A_u^2-bN)}
        [A_u,A_u^2-bN]\me^{-s(A_u^2-bN)}\right]\md s \\
        =& -\str\left[[A_u,A_u^2-bN],\exp\left(-A_u^2+bN\right)\right] \\
        =& b\str\left[[A_u,N],\exp\left(-A_u^2+bN\right)\right] \\
        =& b\str\left[\sqrt{u}(-v+v^*)\exp\left(-A_u^2+bN\right)\right] 
    \end{align*}
    the last equality holds since the following lemma 
    \begin{lemma}
        \begin{equation*}
        [\nabla^E,N] = 0,\quad [v,N] = -v,\quad [v^*,N] = v^*.
        \end{equation*}
    \end{lemma}
    \textcolor{gray}{
        \begin{proof}[Proof of lemma]
        $\nabla^E = \bigoplus \nabla^{E_j}$ preserves the grading of $E$, 
        thus 
        \[
        [\nabla^E,N]=N\nabla^E-N\nabla^E=0.
        \]
        For $u\in C^\infty(X,E_j)$, 
        \begin{equation*}
            [v,N]u = (vN-Nv) = jv u - (j+1)v u = -v u. 
        \end{equation*}
        Similarly, $[v^*,N]u = jv^* u - (j-1)v^* u = v^* u$.
        \end{proof}
    }
    Since $d=\p+\bp$, we can divide the above equation into two parts and 
    set the corresponding parts equal to each other respectively.
    \begin{equation*}
        LHS = \p\str\left[\exp(-A_u^2)\right] + 
        \bp\str\left[\exp(-A_u^2)\right]
    \end{equation*}
    where 
    \begin{gather*}
    \str\left[\exp(-A_u^2+bN)\right] \in 
    \bigoplus_{p,j}\Omega^{p,p}(X,\End(E_j)) \\
    \Longrightarrow
    \left\{
        \begin{aligned}
            &\p\str\left[\exp(-A_u^2+bN)\right] \in 
            \bigoplus_{p,j}\Omega^{p+1,p}(X,\End(E_j)), \\
            &\bp\str\left[\exp(-A_u^2+bN)\right] \in 
            \bigoplus_{p,j}\Omega^{p,p+1}(X,\End(E_j)).
        \end{aligned}
    \right.
    \end{gather*}
    And
    \begin{equation*}
        RHS = -b\sqrt{u}\str\left[v\exp(-A_u^2+bN)\right]
        + b\sqrt{u}\str\left[v^*\exp(-A_u^2+bN)\right]
    \end{equation*}
    where
    \begin{gather*}
        \exp(-A_u^2+bN)\in\bigoplus_{p,q,j}\Omega^{p,q}(X,\Hom(E_j,E_{j+p-q})) \\
        \Longrightarrow
        \left\{
            \begin{aligned}
                &v\exp(-A_u^2+bN) \in 
                \bigoplus_{p,q,j}\Omega^{p,q}(X,\Hom(E_j,E_{j+p-q+1})), \\
                &v^*\exp(-A_u^2+bN) \in 
                \bigoplus_{p,q,j}\Omega^{p,q}(X,\Hom(E_j,E_{j+p-q-1})).
            \end{aligned}
        \right. \\
        \Longrightarrow
        \left\{
            \begin{aligned}
                &\str\left[v\exp(-A_u^2+bN)\right] \in 
                \bigoplus_{p,j}\Omega^{p,p+1}(X,\End(E_j)), \\
                &\str\left[v^*\exp(-A_u^2+bN)\right] \in 
                \bigoplus_{p,j}\Omega^{p+1,p}(X,\End(E_j)).
            \end{aligned}
        \right.
    \end{gather*}
    Thus
    \begin{gather*}
        \left\{
            \begin{aligned}
            &\p\str\left[\exp(-A_u^2+bN)\right] 
            = b\sqrt{u}\str\left[v^*\exp(-A_u^2+bN)\right], \\
            &\bp\str\left[\exp(-A_u^2+bN)\right] 
            = -b\sqrt{u}\str\left[v\exp(-A_u^2+bN)\right].
            \end{aligned}
        \right. \\
        \Longrightarrow\;
        (\p-\bp)\str\left[\exp(-A_u^2+bN)\right]
        = b\str\left[\sqrt{u}V\exp(-A_u^2+bN)\right].
    \end{gather*}
    The last step is to take $\pd{}{b}|_{b=0}$ on both sides, 
    \begin{align*}
        LHS &= (\p-\bp)\pd{}{b}\str\left[\exp(-A_u^2+bN)\right]\Big|_{b=0} \\
        &= (\p-\bp)\int_0^1\str\left[\me^{-(1-s)(A_u^2-bN)}
        \big(-\pd{}{b}(A_u^2-bN)\big)\me^{-s(A_u^2-bN)} \right]\md s\Big|_{b=0} \\
        &= (\p-\bp)\int_0^1\str\left[\me^{-(1-s)A_u^2}N\me^{-sA_u^2}\right]\md s \\
        &= (\p-\bp)\str\left[N\exp(-A_u^2)\right], \\
        RHS &= \pd{}{b}\left(b\str\left[\sqrt{u}V\exp(-A_u^2+bN)\right]\right)\Big|_{b=0} \\
        &= \sqrt{u}\str\left[V\exp(-A_u^2)\right].            
    \end{align*}
    Thus we get formula \eqref{eq:3.2}. It's easy to deduce \eqref{eq:3.3} 
    from \eqref{eq:3.1} and \eqref{eq:3.2}.
    \end{proof}
